\begin{helper}
Let $\mathcal F$ be a $k$-uniform finite multi-hypergraph, and let
$f\in E(\mathcal F)$.  Suppose every vertex of $f$ has degree $D$, and
suppose every edge different from $f$ that meets $f$ meets it in exactly one
vertex.  Then $f$ has exactly $kD-k$ neighboring edges, that is,
\[
\left|\{g\in E(\mathcal F):g\ne f,\ g\cap f\ne\emptyset\}\right|=kD-k .
\]
\end{helper}
\begin{proof}
By \noderef{UniformHypergraph}, the edge $f$ has exactly $k$ vertices.  For a
fixed vertex $v\in f$, \noderef{HypergraphDegree} says that
$\deg_{\mathcal F}(v)$ is the number of edge labels whose edge contains $v$.
The saturation hypothesis gives $\deg_{\mathcal F}(v)=D$.  Among these $D$
edge labels, one is $f$ itself, so exactly $D-1$ labels $g\ne f$ contain
$v$.  Summing over all $v\in f$ gives exactly $k(D-1)$ incidences
\[
(v,g)\quad\text{with}\quad v\in f,\ g\ne f,\ v\in g .
\]

Now compare these incidences with neighboring edges of $f$.  Each neighboring
edge $g\ne f$ contributes at least one incidence because it meets $f$.
The one-intersection hypothesis says it contributes exactly one incidence.
Thus the incidence-counting map from incidences to neighboring edges is
one-to-one and onto after forgetting the unique vertex of $g\cap f$.  Hence
the number of neighboring edges is the same as the incidence count, namely
$k(D-1)=kD-k$.
\end{proof}
